% GENERATED FILE. Edit canonical lessons, then run npm run generate:books.
% canonical-source-hash: fnv1a64:50b6d4f412e83cb5
% canonical-lessons: MR-C499-garodar, MR-C499-balantpan, MR-C499-tanha, MR-C499-rangne, MR-C499-pajne

\chapter{A New Baby}
\label{ch:mr-a-new-baby}
\hlchaptermodality{499}

\begin{chapteropening}
\textbf{By the end of this chapter:} \emph{I can talk about pregnancy, childbirth and a newborn baby.}

Complete the last lesson of chapter 499: I can talk about pregnancy, childbirth and a newborn baby.
\end{chapteropening}

\section[\texorpdfstring{\mr{गरोदर} (garodar)}{गरोदर (garodar)}]{\mr{गरोदर} (garodar) — pregnant}

\label{lesson:MR-C499-garodar}



\begin{quote}
Before the new one: say the Marathi for age, then the Marathi for a teenager.
\end{quote}

\subsection*{You'll want to know: \mr{गरोदर}}
\textbf{\mr{गरोदर}} — \emph{garodar} — \textquotedblleft{}pregnant\textquotedblright{}.

\begin{quote}
\mr{माझी} \mr{मावशी} \mr{गरोदर} \mr{आहे}.
\end{quote}

\begin{quote}
\emph{mājhī māvśī garodar āhe.}
\end{quote}

\begin{quote}
My aunt is pregnant.
\end{quote}

Unchanging; the formal word is \mr{गर्भवती}.

\begin{cousinweb}
\textbf{\mr{गरोदर}} (\emph{garodar}) is usually traced to two Sanskrit words, \emph{garbha}, \textquotedblleft{}womb\textquotedblright{}, and \emph{udara}, \textquotedblleft{}belly\textquotedblright{}.
\end{cousinweb}

\begin{grammarlens}[title={What you've built}]
One more word for describing what you lived and felt.
\end{grammarlens}

\subsection*{Your turn}
\begin{itemize}
\raggedright
  \item \emph{Say it:} \emph{garodar}
  \item \emph{Say it:} \emph{garodar}, once more
  \item \emph{Say it:} \emph{mājhī māvśī garodar āhe.} — My aunt is pregnant.
  \item \emph{Recall:} say the Marathi for age, then the Marathi for a teenager, then say \emph{garodar} again
\end{itemize}

\begin{tcolorbox}[breakable,colback=teal!4,colframe=teal!35!black,title={Before you move on}]
What does \mr{गरोदर} mean? (\textquotedblleft{}Pregnant\textquotedblright{}.) Say it once more.
\end{tcolorbox}
\section[\texorpdfstring{\mr{बाळंतपण} (bāḷantpaṇ)}{बाळंतपण (bāḷantpaṇ)}]{\mr{बाळंतपण} (bāḷantpaṇ) — childbirth; confinement}

\label{lesson:MR-C499-balantpan}



\begin{quote}
Before the new one: say the Marathi for a teenager, then the Marathi for pregnant.
\end{quote}

\subsection*{You'll want to know: \mr{बाळंतपण}}
\textbf{\mr{बाळंतपण}} — \emph{bāḷantpaṇ} — \textquotedblleft{}childbirth; confinement\textquotedblright{}.

\begin{quote}
\mr{मुलीचं} \mr{बाळंतपण} \mr{आईच्या} \mr{घरी} \mr{झालं}.
\end{quote}

\begin{quote}
\emph{mulīcaṁ bāḷantpaṇ āīcyā gharī jhālaṁ.}
\end{quote}

\begin{quote}
The daughter gave birth at her mother's house.
\end{quote}

Neuter.

\mr{बाळंत}, a woman who has recently given birth, + -\mr{पण}.

\begin{cousinweb}
\textbf{\mr{बाळंतपण}} (\emph{bāḷantpaṇ}) is a Marathi derivation: \mr{बाळंत}, a woman who has recently given birth, with \emph{-paṇ}; the origin of \mr{बाळंत} is not settled.
\end{cousinweb}

\begin{grammarlens}[title={What you've built}]
One more word for describing what you lived and felt.
\end{grammarlens}

\subsection*{Your turn}
\begin{itemize}
\raggedright
  \item \emph{Say it:} \emph{bāḷantpaṇ}
  \item \emph{Say it:} \emph{bāḷantpaṇ}, once more
  \item \emph{Say it:} \emph{mulīcaṁ bāḷantpaṇ āīcyā gharī jhālaṁ.} — The daughter gave birth at her mother's house.
  \item \emph{Recall:} say the Marathi for a teenager, then the Marathi for pregnant, then say \emph{bāḷantpaṇ} again
\end{itemize}

\begin{tcolorbox}[breakable,colback=teal!4,colframe=teal!35!black,title={Before you move on}]
What does \mr{बाळंतपण} mean? (\textquotedblleft{}Childbirth; confinement\textquotedblright{}.) Say it once more.
\end{tcolorbox}
\section[\texorpdfstring{\mr{तान्हा} (tānhā)}{तान्हा (tānhā)}]{\mr{तान्हा} (tānhā) — newborn; infant}

\label{lesson:MR-C499-tanha}



\begin{quote}
Before the new one: say the Marathi for pregnant, then the Marathi for childbirth.
\end{quote}

\subsection*{You'll want to know: \mr{तान्हा}}
\textbf{\mr{तान्हा}} — \emph{tānhā} — \textquotedblleft{}newborn; infant\textquotedblright{}.

\begin{quote}
\mr{तान्हं} \mr{बाळ} \mr{दिवसभर} \mr{झोपतं}.
\end{quote}

\begin{quote}
\emph{tānhaṁ bāḷ divasbhar jhoptaṁ.}
\end{quote}

\begin{quote}
A newborn baby sleeps all day.
\end{quote}

Agrees: \mr{तान्हं} \mr{बाळ}, a newborn baby.

\begin{cousinweb}
\textbf{\mr{तान्हा}} (\emph{tānhā}) is inherited: CDIAL lists Marathi \emph{tānhā}, \textquotedblleft{}sucking, of a baby or young animal\textquotedblright{}, under Sanskrit \emph{tṛṣṇaka}, \textquotedblleft{}thirsty\textquotedblright{} (CDIAL 5940).
\end{cousinweb}

\begin{grammarlens}[title={What you've built}]
One more word for describing what you lived and felt.
\end{grammarlens}

\subsection*{Your turn}
\begin{itemize}
\raggedright
  \item \emph{Say it:} \emph{tānhā}
  \item \emph{Say it:} \emph{tānhā}, once more
  \item \emph{Say it:} \emph{tānhaṁ bāḷ divasbhar jhoptaṁ.} — A newborn baby sleeps all day.
  \item \emph{Recall:} say the Marathi for pregnant, then the Marathi for childbirth, then say \emph{tānhā} again
\end{itemize}

\begin{tcolorbox}[breakable,colback=teal!4,colframe=teal!35!black,title={Before you move on}]
What does \mr{तान्हा} mean? (\textquotedblleft{}Newborn; infant\textquotedblright{}.) Say it once more.
\end{tcolorbox}
\section[\texorpdfstring{\mr{रांगणे} (rāṅgṇe)}{रांगणे (rāṅgṇe)}]{\mr{रांगणे} (rāṅgṇe) — to crawl}

\label{lesson:MR-C499-rangne}



\begin{quote}
Before the new one: say the Marathi for childbirth, then the Marathi for newborn.
\end{quote}

\subsection*{You'll want to know: \mr{रांगणे}}
\textbf{\mr{रांगणे}} — \emph{rāṅgṇe} — \textquotedblleft{}to crawl\textquotedblright{}.

\begin{quote}
\mr{आमचं} \mr{बाळ} \mr{आता} \mr{रांगायला} \mr{लागलं}.
\end{quote}

\begin{quote}
\emph{āmcaṁ bāḷ ātā rāṅgāylā lāglaṁ.}
\end{quote}

\begin{quote}
Our baby has started crawling now.
\end{quote}

Past (he …): \mr{तो} \mr{रांगला}, \emph{to rāṅglā}.

\mr{रांगायला} \mr{लागणे}, to start crawling.

\begin{cousinweb}
\textbf{\mr{रांगणे}} (\emph{rāṅgṇe}) is a Marathi verb whose origin is not settled.
\end{cousinweb}

\begin{grammarlens}[title={What you've built}]
One more word for describing what you lived and felt.
\end{grammarlens}

\subsection*{Your turn}
\begin{itemize}
\raggedright
  \item \emph{Say it:} \emph{rāṅgṇe}
  \item \emph{Say it:} \emph{rāṅgṇe}, once more
  \item \emph{Say it:} \emph{āmcaṁ bāḷ ātā rāṅgāylā lāglaṁ.} — Our baby has started crawling now.
  \item \emph{Recall:} say the Marathi for childbirth, then the Marathi for newborn, then say \emph{rāṅgṇe} again
\end{itemize}

\begin{tcolorbox}[breakable,colback=teal!4,colframe=teal!35!black,title={Before you move on}]
What does \mr{रांगणे} mean? (\textquotedblleft{}To crawl\textquotedblright{}.) Say it once more.
\end{tcolorbox}
\section[\texorpdfstring{\mr{पाजणे} (pājṇe)}{पाजणे (pājṇe)}]{\mr{पाजणे} (pājṇe) — to give to drink}

\label{lesson:MR-C499-pajne}



\begin{quote}
Before the new one: say the Marathi for newborn, then the Marathi for to crawl.
\end{quote}

\subsection*{You'll want to know: \mr{पाजणे}}
\textbf{\mr{पाजणे}} — \emph{pājṇe} — \textquotedblleft{}to give to drink\textquotedblright{}.

\begin{quote}
\mr{आईने} \mr{बाळाला} \mr{दूध} \mr{पाजलं}.
\end{quote}

\begin{quote}
\emph{āīne bāḷālā dūdh pājlaṁ.}
\end{quote}

\begin{quote}
The mother fed the baby milk.
\end{quote}

Past, with -\mr{ने} (he …): \mr{त्याने} \mr{बाळाला} \mr{दूध} \mr{पाजलं}, \emph{tyāne bāḷālā dūdh pājlaṁ}.

The one given the drink takes -\mr{ला}: \mr{बाळाला} \mr{पाजणे}.

\begin{cousinweb}
\textbf{\mr{पाजणे}} (\emph{pājṇe}) is a tadbhava, inherited from Sanskrit \emph{pāyayati}, \textquotedblleft{}gives to drink\textquotedblright{}; CDIAL lists Marathi \emph{pājṇẽ} (CDIAL 8098).
\end{cousinweb}

\begin{grammarlens}[title={What you've built}]
One more word for describing what you lived and felt.
\end{grammarlens}

\subsection*{Your turn}
\begin{itemize}
\raggedright
  \item \emph{Say it:} \emph{pājṇe}
  \item \emph{Say it:} \emph{pājṇe}, once more
  \item \emph{Say it:} \emph{āīne bāḷālā dūdh pājlaṁ.} — The mother fed the baby milk.
  \item \emph{Recall:} say the Marathi for newborn, then the Marathi for to crawl, then say \emph{pājṇe} again
\end{itemize}

\begin{tcolorbox}[breakable,colback=teal!4,colframe=teal!35!black,title={Before you move on}]
What does \mr{पाजणे} mean? (\textquotedblleft{}To give to drink\textquotedblright{}.) Say it once more.
\end{tcolorbox}
